\section{GlyphChain release-gate hardening and reproducibility}
\label{sec:glyphchain-release-gate-hardening}

The exact-release evidence gate has been hardened so that role assignment,
independence, time, input provenance and process isolation are machine-enforced
rather than inferred from loosely typed configuration.  These controls address
the difference between having thirty-four evidence entries and proving that each
entry was validly produced, independently approved where required, and evaluated
for the exact candidate under review.

\subsection{Strict independence semantics}

The trust-root inventory is JSON, where booleans and strings are distinct types.
Independence is now recognised only when

\begin{equation}
  \texttt{independent} \equiv \mathrm{true}
\end{equation}

as the literal JSON boolean.  Values such as \texttt{"true"},
\texttt{"false"}, $1$, non-empty objects and non-empty arrays do not grant
independent authority.  This prevents the host language's general truth-value
rules from turning malformed configuration into an elevated release role.

Signer roles must likewise be supplied as an array.  A valid Ed25519 signature
whose trust-root role field has the wrong type cannot satisfy any required role.
Unknown role labels have no effect because a task accepts only its fixed required
role set.

\subsection{Complete separation-of-duties assignment}

A task can require several roles, while a trusted signer may be authorised for
more than one role.  A greedy ``first signer per role'' selection is insufficient:
it can either reuse one person for multiple duties or reject a valid evidence set
when another distinct assignment exists.

For required roles $R=\{r_1,\ldots,r_n\}$ and verified signers $S$, Tessaris now
constructs candidate edges

\begin{equation}
  E = \{(r,s)\mid r\in R,\ s\in S,\ r\in\operatorname{roles}(s)\}.
\end{equation}

The evaluator performs a bounded backtracking search for an injective assignment
$f:R\rightarrow S$ such that $(r,f(r))\in E$ for every role.  Evidence passes
separation of duties only when

\begin{equation}
  r_i \neq r_j \Longrightarrow f(r_i) \neq f(r_j).
\end{equation}

When independence is required, the selected assignment must additionally contain
at least one signer whose trust root carries the literal independent boolean:

\begin{equation}
  \exists r\in R:\operatorname{independent}(f(r))=\mathrm{true}.
\end{equation}

The search is deliberately small and deterministic: a task has only its fixed
role list and an evidence entry accepts at most eight attestations.  It therefore
proves whether a valid assignment exists without introducing an unbounded matching
surface.

\subsection{Observation-time validation}

Every production evidence body must provide a parseable ISO-8601 observation time
with an explicit timezone.  Naive timestamps are rejected because they cannot be
ordered reliably across operators and devices.  The value is normalised to UTC
for comparison and must not occur after the evaluator's report time:

\begin{equation}
  t_{\mathrm{observed}} \leq t_{\mathrm{evaluation}}.
\end{equation}

Malformed and future timestamps leave the task open.  This prevents a pre-dated
approval from being represented as evidence already observed for the candidate.
If operational clocks differ, operators must correct or reconcile the signed
evidence rather than silently applying an evaluator waiver.

\subsection{Exact input provenance}

The readiness report now commits to both complete inputs:

\begin{align}
  h_E &= \operatorname{SHA256}(\operatorname{canon}(E_{1\ldots n})), \\
  h_T &= \operatorname{SHA256}(\operatorname{canon}(T)),
\end{align}

where $E_{1\ldots n}$ is the ordered evidence-entry sequence and $T$ is the
configured trust-root map.  Every task result additionally records:

\begin{itemize}
  \item the canonical evidence-body SHA-256;
  \item the canonical artifact-manifest SHA-256;
  \item the number of retained artifact commitments; and
  \item the signer identifiers present in the evidence entry.
\end{itemize}

These fields allow an auditor to reproduce the decision and detect substitution
of the evidence inventory, trust-root set, body or artifact manifest.  Private
Legal documents and private signing keys are not copied into the report.

The readiness report also contains a SHA-256 integrity checksum over its decision
fields.  This checksum detects accidental or unacknowledged modification when the
expected report digest is retained separately, but it is not described as an
independent digital signature.  Release authenticity derives from re-evaluating
the underlying domain-separated Ed25519 attestations against the independently
administered trust roots.  This distinction avoids overstating a self-hash as a
trust anchor.

\subsection{Inventory schema enforcement}

The command-line release evaluator accepts only the versioned evidence-inventory
and trust-root schemas.  The former must contain an array named
\texttt{evidence\_entries}; the latter must contain an object named
\texttt{trust\_roots}.  Missing, renamed, incorrectly typed or unsupported schema
versions terminate evaluation before a report is written.

This boundary prevents an empty result caused by a spelling error from being
mistaken for a deliberately empty inventory and prevents a partially migrated
trust file from being interpreted under the wrong policy version.

\subsection{Side-effect-free release tooling}

The release evaluator is a security-control utility and must not initialise the
broader business runtime, geometry registries, desktop services or network state
merely because its policy module is imported.  Runtime package exports are now
loaded lazily.  Existing public exports remain available on first access, while
the release module can be imported and its command-line help displayed without
unrelated application startup, circular-import warnings or environment mutation.

This isolation reduces the release tool's operational surface, makes failure
output attributable to the evaluator itself and allows it to run in a minimal
verification environment.

\subsection{Signed-baseline file safety}

The signed baseline generator is part of the release trust chain and now applies
explicit filesystem controls.  Private attestation keys, tracked public trust
anchors, verification inputs, manifest sources and report outputs reject symbolic
links.  This prevents a release command from following an attacker-controlled link
and reading, replacing or signing unintended bytes.

New private keys and public trust anchors use exclusive creation.  A pre-existing
destination is never silently replaced.  The private key is created with owner-only
permissions and flushed to durable storage before use.  Baseline reports are
written to a uniquely created file in the destination directory, flushed, assigned
the intended public-read mode and atomically renamed over the preceding regular
report.  The containing directory is then synchronised.  A crash therefore leaves
either the previous complete report or the replacement complete report, rather
than a partially written signed artefact.

The source-manifest builder also rejects symbolic links among the explicitly
versioned files and recursively collected GlyphChain sources.  Every manifest
digest consequently represents bytes resident in the intended release tree rather
than a target that can be redirected after review.  The signed report retains the
ordered path-and-SHA--256 entry array as well as its aggregate digest and count, so
an auditor can identify the exact files covered rather than trusting an opaque
aggregate alone.  The core readiness evaluator, evidence and trust-root tools,
baseline generator, regression tests and every GlyphChain Markdown and \LaTeX{}
technical document are mandatory members.  A regression compares that mandatory
set with the generated source inventory and fails if a release-decision component
or a technical document is omitted.

\subsection{Regression evidence and claim boundary}

The focused regression set proves strict boolean independence, non-greedy distinct
role assignment, malformed and future-time rejection, exact evidence/trust input
binding, checklist-to-register equality and side-effect-free command-line import.
Baseline regressions additionally prove private-key, public-key and report symlink
refusal, preservation of the unintended target, atomic replacement and final file
permissions.
The broader retained GlyphChain and Legal suite is rerun before the signed baseline
is replaced.

These controls establish that the local evaluator can make a reproducible,
fail-closed decision from supplied signed evidence.  They do not create the live
provider events, independent reviews, qualified Legal approvals, multi-device
observations or customer-pilot outcomes needed to satisfy the thirty-four external
release gates.
